\begin{proof}
\textbf{1. Inner‑product expansion.}  
By linearity in the first argument and conjugate symmetry,
\begin{align*}
\langle x+y,\,x+y\rangle
&=\langle x,x\rangle+\langle x,y\rangle+\langle y,x\rangle+\langle y,y\rangle  \\
&=\|x\|^{2}+2\,\operatorname{Re}\langle x,y\rangle+\|y\|^{2}.
\end{align*}
In a real inner‑product space \(\operatorname{Re}\langle x,y\rangle=\langle x,y\rangle\); in the complex case the argument below is unchanged because
\(|\langle x,y\rangle|\ge\operatorname{Re}\langle x,y\rangle\).

\medskip
\textbf{2. Cauchy–Schwarz inequality.}  
\[
|\langle x,y\rangle|\le \|x\|\,\|y\|,
\]
with equality iff \(x\) and \(y\) are linearly dependent.

\medskip
\textbf{3. Bounding the mixed term.}  
Since \(\operatorname{Re}\langle x,y\rangle\le |\langle x,y\rangle|\),
\[
2\,\operatorname{Re}\langle x,y\rangle\le 2\,|\langle x,y\rangle|
\le 2\,\|x\|\,\|y\|.
\]

\medskip
\textbf{4. Inequality for the squared norm.}  
Insert the estimate from step 3 into the expansion of step 1:
\begin{align*}
\|x+y\|^{2}
&=\|x\|^{2}+2\,\operatorname{Re}\langle x,y\rangle+\|y\|^{2} \\
&\le \|x\|^{2}+2\,\|x\|\,\|y\|+\|y\|^{2}
   =\bigl(\|x\|+\|y\|\bigr)^{2}.
\end{align*}

\medskip
\textbf{5. Taking square roots.}  
Norms are non‑negative, hence the square‑root function is monotone on
\([0,\infty)\). From the previous inequality we obtain
\[
\|x+y\|\le \|x\|+\|y\|.
\]

\medskip
\textbf{6. Degenerate cases.}  
If \(x=0\) (or \(y=0\)) the inequality reduces to \(\|y\|\le\|y\|\) (or the symmetric
statement), which is true. When both vectors are zero both sides equal \(0\).

\medskip
\textbf{7. Equality condition.}  
Equality in step 4 requires equality in both inequalities of step 3:

\begin{enumerate}[(a)]
\item \(\operatorname{Re}\langle x,y\rangle = |\langle x,y\rangle|\) forces
      \(\langle x,y\rangle\) to be a non‑negative real number.
\item \(|\langle x,y\rangle| = \|x\|\,\|y\|\) forces linear dependence, i.e.
      there exists \(\alpha\in\mathbb{C}\) with \(y=\alpha x\).
\end{enumerate}
Writing \(y=\alpha x\) we have \(\langle x,y\rangle=\alpha\|x\|^{2}\).
Condition (a) then gives \(\alpha\|x\|^{2}\in[0,\infty)\), so \(\alpha\) is a
non‑negative real scalar. Setting \(\lambda:=\alpha\ge 0\) yields \(y=\lambda x\).

Conversely, if \(y=\lambda x\) with \(\lambda\ge 0\), then
\[
\|x+y\|=\|(1+\lambda)x\|=(1+\lambda)\|x\|
       =\|x\|+\|\lambda x\|=\|x\|+\|y\|,
\]
so equality holds. The cases \(x=0\) or \(y=0\) correspond to \(\lambda=0\).

Hence equality occurs exactly when the two vectors are positively collinear,
i.e. there exists \(\lambda\ge 0\) such that \(y=\lambda x\). \(\square\)
\end{proof}